\section{Experiment}
\label{sec:experiments}

We evaluate \system on three open-source persistent-memory systems:
Mem0~\citep{chhikara2025mem0}, A-Mem~\citep{xu2025amem}, and
Graphiti~\citep{rasmussen2025zep}. We use
LoCoMo~\citep{maharana2024evaluating} and
LongMemEval-S~\citep{wu2024longmemeval} as benchmark seed corpora. Results are
reported separately by benchmark and backend, with repeated campaigns
aggregated over repetition seeds.

\paragraph{Metrics.}
UF@B counts distinct canonical semantic fault signatures among evaluator-
confirmed failures in the first \(B\) valid fuzzing executions. The CFS schema
and equivalence rule are fixed; executable canonicalization and failure-surface
assignment depend on the separately frozen normalization and evaluator
predicates. Cov@B is the fraction of the frozen initialized retrievable memory
entries reached during those \(B\) executions. Coverage starts empty and
initialization retrievals receive no credit. A backend--benchmark condition
with an empty campaign-level inventory is RQ1-ineligible: Cov@B is undefined
and UF@B is not reported. An individual empty checkpoint does not exclude a
campaign whose campaign-level inventory is non-empty.

\subsection{RQ1: How effective is U-Fuzz at memory-use failure discovery?}
\label{subsec:rq1}

We compare full \system with four same-space baselines: Random Mutation,
Unguided LLM, LLM-as-Judge, and Coverage-Guided. These methods share the
benchmark seed corpus, eligible checkpoint set, backend initialization
protocol, full applicable mutation space, validation procedure, generation-
attempt limits, execution budget, and seed-retention eligibility. We also
evaluate U-Fuzz-Q and U-Fuzz-M as restricted query-only and memory-only
ablations.

\begin{table}[t]
\centering
\caption{
Fuzzing effectiveness under a fixed execution budget \(B\).
UF@\(B\) is the number of unique confirmed faults, and Cov@\(B\)
is memory-entry coverage.
}
\label{tab:rq1}

\renewcommand{\arraystretch}{1.02}
\setlength{\tabcolsep}{3.5pt}

\begin{minipage}{0.8\linewidth}
\centering
\begin{tabularx}{\linewidth}{
    >{\raggedright\arraybackslash}p{0.25\linewidth}
    *{6}{>{\centering\arraybackslash}X}
}
\toprule
\textbf{Method}
& \multicolumn{2}{c}{\textbf{Mem0}}
& \multicolumn{2}{c}{\textbf{A-Mem}}
& \multicolumn{2}{c}{\textbf{Graphiti}} \\
\cmidrule(lr){2-3}
\cmidrule(lr){4-5}
\cmidrule(lr){6-7}
& {\small \textbf{UF@\(B\)}} & {\small \textbf{Cov@\(B\)}}
& {\small \textbf{UF@\(B\)}} & {\small \textbf{Cov@\(B\)}}
& {\small \textbf{UF@\(B\)}} & {\small \textbf{Cov@\(B\)}} \\
\midrule

\multicolumn{7}{l}{\textbf{LoCoMo}} \\[1pt]

\rowcolor[HTML]{F3E6D5}
Random Mutation
& -- & -- & -- & -- & -- & -- \\
\rowcolor[HTML]{F3E6D5}
Unguided LLM
& -- & -- & -- & -- & -- & -- \\
\rowcolor[HTML]{F3E6D5}
LLM-as-Judge
& -- & -- & -- & -- & -- & -- \\
\rowcolor[HTML]{F3E6D5}
Coverage-Guided
& -- & -- & -- & -- & -- & -- \\

\noalign{\vskip 1pt}

\rowcolor[HTML]{FFF9F2}
U-Fuzz-Q
& -- & -- & -- & -- & -- & -- \\
\rowcolor[HTML]{FFF9F2}
U-Fuzz-M
& -- & -- & -- & -- & -- & -- \\
\rowcolor[HTML]{FFF9F2}
\textbf{U-Fuzz}
& -- & -- & -- & -- & -- & -- \\

\addlinespace[3pt]
\multicolumn{7}{l}{\textbf{LongMemEval-S}} \\[1pt]

\rowcolor[HTML]{F3E6D5}
Random Mutation
& -- & -- & -- & -- & -- & -- \\
\rowcolor[HTML]{F3E6D5}
Unguided LLM
& -- & -- & -- & -- & -- & -- \\
\rowcolor[HTML]{F3E6D5}
LLM-as-Judge
& -- & -- & -- & -- & -- & -- \\
\rowcolor[HTML]{F3E6D5}
Coverage-Guided
& -- & -- & -- & -- & -- & -- \\

\noalign{\vskip 1pt}

\rowcolor[HTML]{FFF9F2}
U-Fuzz-Q
& -- & -- & -- & -- & -- & -- \\
\rowcolor[HTML]{FFF9F2}
U-Fuzz-M
& -- & -- & -- & -- & -- & -- \\
\rowcolor[HTML]{FFF9F2}
\textbf{U-Fuzz}
& -- & -- & -- & -- & -- & -- \\

\bottomrule
\end{tabularx}
\end{minipage}
\end{table}

Table~\ref{tab:rq1} reports UF@\(B\) and Cov@\(B\) for each benchmark and
memory system. On LoCoMo, \system finds \textbf{[XXX]}, \textbf{[XXX]}, and
\textbf{[XXX]} unique faults on Mem0, A-Mem, and Graphiti, with corresponding
coverage of \textbf{[XXX]}, \textbf{[XXX]}, and \textbf{[XXX]}. On
LongMemEval-S, the corresponding results are \textbf{[XXX]}, \textbf{[XXX]},
and \textbf{[XXX]} for UF@\(B\), and \textbf{[XXX]}, \textbf{[XXX]}, and
\textbf{[XXX]} for Cov@\(B\). The comparison with U-Fuzz-Q and U-Fuzz-M
further shows how query and memory-state mutations contribute to the full
fuzzing process.
